% Research notes for the parent article.  Not a standalone document.
% Source snapshot: 3447bc59b78d138a7f0d536b66ac6e5cc5d5e49f.

\section{Integral diagonalization of an odd-weight shuffle minor}
\label{oi:sec:shuffle}

The following is a theorem about an integral relation presentation.
It says nothing about arithmetic independence of evaluated periods.
Its immediate analytic consequence is a uniformly justified choice of
Gaussian double-polylogarithm coordinates at every odd weight.

\subsection{The relation matrix and its arithmetic structure}

Let $m\geq1$, set $w=2m+1$, and use the convention
\[
 F_{a,b}(z)=\operatorname{Li}_{a,b}(z,1)
 =\sum_{n>r\geq1}\frac{z^n}{n^a r^b}.
\]
The standard same-point integral shuffle gives
\begin{equation}\label{oi:eq:shuffle}
 \operatorname{Li}_{p}(z)\operatorname{Li}_{w-p}(z)
 =\sum_{a=1}^{2m}
 \left\{\binom{a-1}{p-1}+\binom{a-1}{2m-p}\right\}F_{a,w-a}(z),
 \qquad 1\leq p\leq m.
\end{equation}
Binomial coefficients outside their usual range are zero.  This holds
as a holomorphic identity for $|z|<1$, and also at $z=i$ with the
ordinary convergent boundary values.  In particular, the possible outer
index one is harmless there: summation by parts applied to
$H_{n-1}^{(b)}/n$ gives ordinary convergence, and the corresponding
iterated integral has no singular endpoint letter.

Select the columns with odd outer index $a=2k+1$, so their inner indices
are even.  With $p,k$ now zero based, the resulting $m\times m$ matrix is
\begin{equation}\label{oi:eq:M}
 (M_m)_{p,k}=
 \binom{2k}{p}+\binom{2k}{2m-1-p},
 \qquad 0\leq p,k<m.
\end{equation}

\begin{theorem}[Integral diagonalization of the shuffle minor]
\label{oi:thm:integral}
There are explicitly constructible matrices $U_m,V_m\in
\operatorname{GL}_m(\mathbb Z)$ such that
\begin{equation}\label{oi:eq:diagonal}
 U_mM_mV_m=\operatorname{diag}(1,3,5,\ldots,2m-1).
\end{equation}
Consequently:
\begin{enumerate}
 \item $\det M_m=(2m-1)!!$.
 \item $\operatorname{coker}_{\mathbb Z}M_m\cong
       \bigoplus_{j=0}^{m-1}\mathbb Z/(2j+1)\mathbb Z$.
       This direct-sum display is an integral diagonal form, not generally
       the ordered Smith invariant-factor display.
 \item The least positive integer $L$ for which every entry of
       $LM_m^{-1}$ is an integer is
       $L_m=\operatorname{lcm}(1,3,5,\ldots,2m-1)$.
 \item Over characteristic two the matrix has full rank.  Over a field
       of odd characteristic $\ell$ its rank is
       \[
       m-\left\lfloor\frac{2m+\ell-1}{2\ell}\right\rfloor.
       \]
\end{enumerate}
\end{theorem}

\begin{proof}
Put $n=2m-1$.  The palindromic polynomials
\[
 Q_k(x)=(1+x)^{2k}(1+x^{n-2k}),\qquad 0\leq k<m,
\]
have degree $n$, and their coefficients of $x^p$, $0\leq p<m$, are
precisely $(M_m)_{p,k}$.  Define
\[
 b_j(x)=x^j(1+x)^{n-2j},\qquad 0\leq j<m,
 \qquad
 B_{p,j}=\binom{n-2j}{p-j}.
\]
The coefficient matrix $B$ is lower triangular with diagonal entries
one, hence belongs to $\operatorname{GL}_m(\mathbb Z)$.

Let $E_r(t)$ be the polynomial defined by
\[
 E_r(y(1-y))=y^r+(1-y)^r.
\]
Its existence with integer coefficients follows, for example, from
$E_0=2$, $E_1=1$, and $E_r=E_{r-1}-tE_{r-2}$.
Let $E$ be the $m\times m$ integer matrix whose column $j$ is the
coefficient vector of $E_{2j+1}(t)$ in $1,t,\ldots,t^{m-1}$.
Under $y=x/(1+x)$, $t=x/(1+x)^2$, one has
\[
 Q_k(x)=(1+x)^n E_{n-2k}\!\left(\frac{x}{(1+x)^2}\right).
\]
If $R$ is the column-reversal matrix, this identity gives the exact
factorization
\begin{equation}\label{oi:eq:BER}
 M_m=BER.
\end{equation}

It remains to diagonalize $E$ over the integers.  Use the monic odd
polynomials
\[
 p_j(y)=y\prod_{r=1}^j(y^2-r^2),\qquad 0\leq j<m,
\]
and the integer polynomials
\[
 q_j(t)=(-1)^j\prod_{r=1}^j\bigl(t+r(r-1)\bigr).
\]
Their respective leading coefficients are $1$ and $(-1)^j$.
Let $P$ be the coefficient matrix of $p_j$ in
$y,y^3,\ldots,y^{2m-1}$, and let $Q$ be the coefficient matrix of
$q_j$ in $1,t,\ldots,t^{m-1}$.  Both are integer triangular matrices
with unit diagonal up to signs, so both are unimodular.

The central-factorial finite-difference identity is
\begin{align*}
 p_j(y)+p_j(1-y)
 &=p_j(y)-p_j(y-1)\\
 &=(2j+1)\prod_{r=1-j}^{j}(y-r)\\
 &=(2j+1)(-1)^j
       \prod_{r=1}^j\bigl(y(1-y)+r(r-1)\bigr)\\
 &=(2j+1)q_j(y(1-y)).
\end{align*}
For the middle equality, write
$p_j(y)=(y+j)^{\underline{2j+1}}$ and use
$z^{\underline n}-(z-1)^{\underline n}
=n(z-1)^{\underline{n-1}}$; this is a polynomial identity and needs no
restriction on $y$.  Since substitution $t\mapsto y(1-y)$ is injective
on $\mathbb Z[t]$, comparison in the $t$ basis gives
\[
 EP=QD,\qquad D=\operatorname{diag}(1,3,\ldots,2m-1).
\]
Combining this with \eqref{oi:eq:BER},
\[
 U_m=(BQ)^{-1},\qquad V_m=RP
\]
are integral unimodular matrices and satisfy \eqref{oi:eq:diagonal}.

The determinant sign is also fixed by these formulas:
$\det B=\det P=1$, while
$\det Q=\det R=(-1)^{m(m-1)/2}$.  Thus $\det M_m=\det D>0$.
The cokernel and field-rank assertions follow from integral equivalence.
For odd $\ell$, the number of positive odd multiples of $\ell$ not
exceeding $2m-1$ is the displayed floor.

Finally, $M_m^{-1}=V_mD^{-1}U_m$, with integral $U_m,V_m$ and integral
inverses.  Therefore $LM_m^{-1}$ is integral if and only if $LD^{-1}$
is integral, if and only if every $2j+1$ divides $L$.  This proves the
optimal denominator assertion.
\end{proof}

\paragraph{Smith factors without ambiguity.}
For each prime $\ell$, sort the $m$ integers
$v_\ell(1),v_\ell(3),\ldots,v_\ell(2m-1)$ into
$e_{\ell,1}\leq\cdots\leq e_{\ell,m}$.  The ordered Smith invariant
factors are exactly
$d_r=\prod_\ell\ell^{e_{\ell,r}}$.  For $m=5$, for example, they are
$(1,1,1,3,315)$, while the determinant is $945$ and the optimal inverse
denominator is $315$.  At $m=8$ they are
$(1,1,1,1,1,3,15,45045)$.  Thus determinant size and the required common
denominator of the inverse are quite different quantities.

\subsection{The uniform odd-inner-index basis}

\begin{corollary}[An integral justification for the discovery basket]
\label{oi:cor:basket}
In the rational vector space on the $2m$ formal double coordinates
$F_{a,2m+1-a}$, modulo the $m$ same-point shuffle product rows, the $m$
classes
\[
 F_{2m,1},F_{2m-2,3},\ldots,F_{2,2m-1}
\]
form a basis.  Before taking the quotient, every double is an explicit
linear combination of these odd-inner-index coordinates and products
$\operatorname{Li}_p\operatorname{Li}_{2m+1-p}$, with all coefficients
in $L_m^{-1}\mathbb Z$.
\end{corollary}
\begin{proof}
Separate \eqref{oi:eq:shuffle} into even-inner coordinates $X$ and
odd-inner coordinates $Y$.  If $N$ is the integer coefficient matrix
of the latter, the equation is $M_mX+NY=P$, where $P$ is the product
vector.  Hence $X=M_m^{-1}P-M_m^{-1}NY$.  The theorem proves both
existence and the coefficient-ring assertion.  It also proves that the
$m$ rows have rank $m$, so the $m$ surviving coordinates are independent
in this specified formal quotient.  Additional functional equations
or numerical evaluation may introduce further relations.
\end{proof}

\paragraph{Relation to Zagier's matrix.}
Zagier's 2012 paper \cite{Zagier}, Section 5, equation (36), defines the
$(m-1)\times(m-1)$ matrix
\[
 (A_m)_{r,s}=\binom{2m-2s}{2r-1}
             +\binom{2m-2s}{2m-2r},\qquad 1\leq r,s<m.
\]
It is permutation-equivalent to the minor of $M_m$ obtained by removing
row zero and column zero: reverse the columns, and fold the odd row
indices $2r-1$ using $p\leftrightarrow2m-1-p$.  These folded row
indices cover $1,\ldots,m-1$ exactly once.  The first column of $M_m$
is $(1,0,\ldots,0)^T$, so integer column subtraction removes its first
row's other entries.  Therefore
\[
 A_m\ \text{is integrally equivalent to}\
 \operatorname{diag}(3,5,\ldots,2m-1).
\]
In particular $|\det A_m|=(2m-1)!!$, the determinant formula conjectured
in Zagier's equation (37); its sign is determined by the explicit
permutations.  The paper's inverse formula was proved by Ding Ma~\cite{Ma} and subsequently by Yawen Ma and
Lee-Peng Teo~\cite{MaTeo},
\emph{Another Proof of Zagier's Matrix Conjecture}, Journal of Integer
Sequences 25 (2022), Article 22.6.4.  We do not claim priority for the
determinant identity or for the known inverse formula.  The direct
integral diagonalization above supplies all Smith factors and the exact
denominator statement in the manuscript's shuffle coordinates.

\subsection{A short exact weight-eleven identity}

The arithmetic theorem supplies concrete identities as well as a basis
count.  A useful example has a four-row certificate.

\begin{corollary}[A weight-eleven Gaussian double reduction]
\begin{align*}
 &105g_{9,2}+210g_{10,1}+35g_{8,3}-7g_{6,5}+5g_{4,7}-7g_{2,9}\\
 &\quad=\frac{223339}{29727129600}\pi^{11}
 +\frac{1785}{131072}\beta(2)\zeta(9)
 +\frac{63}{512}\beta(4)\zeta(7)\\
 &\qquad\quad-\frac{15}{64}\beta(6)\zeta(5)
 -\frac{21}{16}\beta(8)\zeta(3).
\end{align*}
\end{corollary}
\begin{proof}
At general $z$ the stronger functional identity is
\begin{align*}
 &105F_{9,2}(z)+210F_{10,1}(z)+35F_{8,3}(z)
       -7F_{6,5}(z)+5F_{4,7}(z)-7F_{2,9}(z)\\
 &\quad=-7\operatorname{Li}_2(z)\operatorname{Li}_9(z)
       +14\operatorname{Li}_3(z)\operatorname{Li}_8(z)
       -16\operatorname{Li}_4(z)\operatorname{Li}_7(z)
       +8\operatorname{Li}_5(z)\operatorname{Li}_6(z).
\end{align*}
It is $-7,14,-16,8$ times the shuffle rows with $p=2,3,4,5$,
respectively.  Direct coefficient comparison gives, in increasing
outer-index order, the vector
$(0,-7,0,5,0,-7,0,35,105,210)$.  This is the complete small
certificate of the functional identity.  Specialize to $z=i$, take
imaginary parts, and use
\[
 \operatorname{Re}\operatorname{Li}_k(i)
 =-2^{-k}(1-2^{1-k})\zeta(k),\qquad k>1,
 \quad
 \operatorname{Im}\operatorname{Li}_k(i)=\beta(k).
\]
The even-zeta and odd-beta evaluations give the displayed rational
coefficient of $\pi^{11}$ and the remaining four products.
\end{proof}

\section{Frozen Gaussian candidates at weights eleven and thirteen}
\label{sec:candidates}

Retain $g_{a,b}=\operatorname{Im}F_{a,b}(i)$ and
$S_p=\sum_{n\geq0}(-1)^nH_n/(2n+1)^p$.
The following identity is conjectural; the exact matrix theorem above
justifies the choice of comparison coordinates, but does not imply
membership of the mixed sum in their span.

\begin{conjecture}[Weight eleven]\label{conj:S10}
\begin{align*}
 S_{10}={}&\frac{570649054}{415912609}g_{10,1}
 +\frac{30750200}{415912609}g_{8,3}
 -\frac{14642272}{415912609}g_{6,5}\\
 &+\frac{15009920}{415912609}g_{4,7}
 -\frac{641796160}{1247737827}g_{2,9}
 +\frac{15870231535729}{2649404577860812800}\pi^{11}\\
 &-\frac{852385525}{851789023232}\beta(2)\zeta(9)
 -\frac{88141815}{7605259136}\beta(4)\zeta(7)\\
 &-\frac{437907815}{6654601744}\beta(6)\zeta(5)
 -\frac{468343007}{1663650436}\beta(8)\zeta(3)
 -2\beta(10)\log2.
\end{align*}
\end{conjecture}

In the ordered basket
\[
 (S_{10},g_{10,1},g_{8,3},g_{6,5},g_{4,7},g_{2,9},\pi^{11},
 \beta(2)\zeta(9),\beta(4)\zeta(7),\beta(6)\zeta(5),
 \beta(8)\zeta(3),\beta(10)\log2),
\]
the frozen primitive integer vector is
\[
\begin{split}
 (&2649404577860812800,-3635091082366156800,
 -195881824419840000,93272725153382400,\\
 &-95614679384064000,1362768401793024000,-15870231535729,
 2651259936960000,\\
 &30705505754112000,174344763875328000,
 745847853554073600,5298809155721625600).
\end{split}
\]
The search used 1600 finite Euler weights computed as exact integers,
550-digit arithmetic for the values, then PSLQ at 300 working digits,
tolerance $10^{-280}$, coefficient bound $10^{30}$, and at most 30000
iterations.  After freezing the vector, the residual evaluated using the
retained 550-digit values was approximately $2.575\times10^{-464}$.
This number is a diagnostic and is consistent with the finite Euler
truncation.  The accompanying exact rational certificate is the
authoritative finite proximity claim.  Neither small residuals nor a
residual interval containing zero prove this conjecture.

An independent 120-digit quadrature audit used
\[
 S_{10}=-\frac1{9!}\int_0^\infty
 \frac{t^9e^{-t}\log(1+e^{-2t})}{1+e^{-2t}}\,dt,
\]
and
\[
 g_{a,b}=\frac1{(a-1)!}\int_0^1(-\log x)^{a-1}
 \operatorname{Im}\frac{i\operatorname{Li}_b(ix)}{1-ix}\,dx.
\]
For $b=1$, it instead used the integrated logarithmic formula
\[
 g_{a,1}=-\frac1{2(a-2)!}\int_0^1
 \frac{(-\log x)^{a-2}\log(1+x^2)\arctan x}{x}\,dx.
\]
The normalized residual was approximately $3.49\times10^{-121}$.
This confirms the normalization by a different representation, with
floating-point quadrature rather than the finite Euler sums.  No
rigorous quadrature error estimate is asserted, and this audit is not
a premise of the rational proximity proof.

\begin{proposition}[Certified weight-eleven proximity]\label{prop:S10-proximity}
Let $R_{10}$ be $S_{10}$ minus the displayed conjectural right side.
Then
\[
 -10^{-650}<R_{10}<10^{-650}.
\]
The delivered rational residual interval contains zero.  This statement
is a proved enclosure for the normalized residual; the conjectural
assertion $R_{10}=0$ remains open.
\end{proposition}
\begin{proof}
Apply the rational bounds proved immediately below with $N=2200$ Euler
terms, 600 terms in each of the two arctangent series, and 800 terms in
the logarithm series.  The standard-library program
\texttt{verify\_gaussian\_proximity.py} reconstructs all twelve basket
intervals from these finite integer/rational formulas, reads the frozen
primitive vector from \texttt{gaussian\_candidates.json}, and performs
ordinary exact rational interval addition and multiplication.  It
divides the integer dot product by its first coefficient
$2649404577860812800$, so the output is precisely $R_{10}$, rather than
the unnormalized integer residual.  Direct rational comparison proves
that the lower endpoint exceeds $-10^{-650}$ and the upper endpoint is
below $10^{-650}$.  Exact endpoints of this residual and all basket
values are in \texttt{gaussian\_proximity\_certificate.json}.
No floating-point approximation enters this verification.
\end{proof}

\subsection{A further candidate at weight thirteen}

The same structured basket yields one additional frozen conjecture.
Let
\[
\begin{split}
 \mathbf b_{12}={}&(S_{12},g_{12,1},g_{10,3},g_{8,5},g_{6,7},g_{4,9},
 g_{2,11},\pi^{13},\\
 &\beta(2)\zeta(11),\beta(4)\zeta(9),\beta(6)\zeta(7),
 \beta(8)\zeta(5),\beta(10)\zeta(3),\beta(12)\log2).
\end{split}
\]
The conjecture is $\mathbf c_{12}\cdot\mathbf b_{12}=0$, where
\[
\begin{split}
\mathbf c_{12}=(&67206362357091388096512000,
-92814968433857420014387200,\\
&-5035645608385269517516800,
2454053342295358816911360,\\
&-2557982490705884558131200,
4600624150682630105333760,\\
&10639593135931482596966400,-43411484181262303199,\\
&5190040482548669193600,71048339137356985958400,\\
&743250995585205475737600,5397522019420435985203200,\\
&22801241978065356167577600,134412724714182776193024000).
\end{split}
\]
It has greatest common divisor one.  A 440-digit PSLQ calculation with
tolerance $10^{-420}$, coefficient bound $10^{30}$, and at most 50000
iterations produced this vector from the 550-digit/1600-Euler-term
values.  Its retained-precision integer residual was approximately
$2.581\times10^{-457}$.
The same exact rational calculation as for $S_{10}$ proves
\[
 \left|\frac{\mathbf c_{12}\cdot\mathbf b_{12}}
 {67206362357091388096512000}\right|<10^{-650}.
\]
The delivered normalized interval contains zero.  This remains a
conjectural identity, supported by a proved proximity enclosure.

The first, weaker search at 300 digits and tolerance $10^{-280}$
returned a different vector.  Its higher-precision residual failed
immediately.  That discarded vector is preserved in
\texttt{gaussian\_candidates.json}, separately from the retained vector.
Using the same rigorous basket intervals, its \emph{unnormalized}
integer residual $R_{\mathrm{bad}}$ satisfies
\[
 -\frac{2191}{10^{277}}<R_{\mathrm{bad}}
 <-\frac{2190}{10^{277}}<0.
\]
Thus the original weight-thirteen vector is rigorously false, not merely
unconfirmed.  The exact endpoints are included in the certificate and
the optional \texttt{--s12} replay verifies both the retained proximity
statement and this rejection.  This is a practical example of why
discovery precision must not be reused as validation precision.

\subsection{The exact rational error bound used by the certificate}

For clarity, the error bound is rederived here.  It is the direct
one-sided Euler bound already used in the pinned manuscript, not a
new proof of the proposed identity.  For a sequence $f_n$, write
\[
 \Delta^kf_0=\sum_{n=0}^k(-1)^n\binom{k}{n}f_n,
 \qquad
 E_N(f)=\sum_{k=0}^{N-1}\frac{\Delta^kf_0}{2^{k+1}}.
\]
The exact finite identity
\begin{equation}\label{oi:eq:euler-finite}
 E_N(f)=\frac1{2^N}\sum_{n=0}^{N-1}(-1)^n
 \left(\sum_{r=n+1}^{N}\binom Nr\right)f_n
\end{equation}
follows by interchanging the two finite sums and using Pascal's
identity.  It is convenient computationally because its weights are
integers divided by one power of two; the implementation does not
form an unstable floating-point finite-difference triangle.

For $a\geq2$ and integer $b\geq1$, put
$f_n=H_{2n}^{(b)}/(2n+1)^a$.  The elementary Mellin formulas give
\[
 f_n=\frac1{\Gamma(a)\Gamma(b)}
 \int_0^1\int_0^1
 \frac{(-\log x)^{a-1}(-\log t)^{b-1}}{1-t}
 \bigl((x^2)^n-(x^2t^2)^n\bigr)\,dt\,dx.
\]
Each grouped difference is integrable, including at $b=1$; the
individual terms should not be separated there.  Finite differences
turn the bracket into
$(1-x^2)^k-(1-x^2t^2)^k$, which is nonpositive when $k\geq1$.
Summing the geometric tail yields the exact remainder expression
\begin{align*}
 g_{a,b}-E_N(f)
 =\frac{2^{-N}}{\Gamma(a)\Gamma(b)}
 \int_0^1\int_0^1
 \frac{(-\log x)^{a-1}(-\log t)^{b-1}}{1-t}
 \left\{\frac{(1-x^2)^N}{1+x^2}
       -\frac{(1-x^2t^2)^N}{1+x^2t^2}\right\}\,dt\,dx.
\end{align*}
One justification is to start with the convergent original series:
for $a\geq2$ it is absolutely summable since $b\geq1$.  The displayed
integral equals its sum by dominated convergence of grouped kernels.
Alternatively, the uniform bound below makes the series of Euler
remainders absolutely integrable and identifies the same sum.

The function $B_N(u)=(1-u)^N/(1+u)$ decreases on $[0,1]$ and obeys
$|B_N'(u)|\leq N+1$.  Thus the bracket lies between
$-(N+1)x^2(1-t^2)$ and zero.  Integrating this bound gives
\begin{equation}\label{oi:eq:gaussian-bound}
 E_N(f)-\frac{N+1}{2^N3^a}\left(1+2^{-b}\right)
 \leq g_{a,b}\leq E_N(f).
\end{equation}
The integral of $x^2(1-t^2)$ against the positive Mellin weight is
exactly $f_1=H_2^{(b)}/3^a$, explaining every factor in the bound.
For $f_n=H_n/(2n+1)^p$, $p\geq2$, replace $x^2t^2$ by $x^2t$.
The same argument gives
\begin{equation}\label{oi:eq:mixed-bound}
 E_N(f)-\frac{N+1}{2^N3^p}\leq S_p\leq E_N(f).
\end{equation}

The single-value sequences $(2n+1)^{-s}$ and $(n+1)^{-s}$ are
Hausdorff moment sequences with mass one.  Their omitted Euler tails
are nonnegative and bounded by $2^{-N}$.  This encloses $\beta(s)$
and $\eta(s)$; the relation
$\zeta(s)=\eta(s)/(1-2^{1-s})$ encloses every odd zeta value needed.
The elementary constants are evaluated independently with
\[
 \pi=16\arctan(1/5)-4\arctan(1/239),\qquad
 \log2=2\sum_{n\geq0}\frac1{(2n+1)3^{2n+1}}.
\]
The alternating arctangent series is enclosed by its first omitted
term.  For $M$ terms in the logarithm series, the positive tail is at
most $9/[4(2M+1)3^{2M+1}]$.  Machin's identity uses the real principal
arctangent: the tangent of $4\arctan(1/5)-\arctan(1/239)$ is one,
and this angle lies in $(0,\pi/2)$, fixing the branch.

All these bounds are rational.  The certificate applies
\eqref{oi:eq:euler-finite} using an integer common denominator
$D=\operatorname{lcm}(1,\ldots,2N)$, and performs only integer and
rational interval operations.  It encloses every basket value before
forming the fixed linear combination.  Thus it does not assume that
the conjectural residual is zero at any stage.

